\section{Discussion}

\subsection{Why Adaptive Solves the K=200 Hazard}

The Paper~7 K=200 hazard arose from forcing the procedure to use
lag-1-fit weights in every window, including windows where the
calibration-target shift was too large for the lag-1 fit to
generalise. The adaptive detector identifies those windows directly
through the magnitude of the calibration-target shift and gates lag-1
out of them. At $K = 200$ this gating fires four of five times --- the
detector and the capacity-conditional reality agree at this capacity.

Paper~8's smoothers attempted the opposite intervention: keep using
fitted weights but draw them from a longer history. That made the
problem worse because longer history brings in more mis-aligned
calibration targets. The adaptive detector instead acknowledges
when the fit is mis-aligned and falls back to a known-safe baseline.
The asymmetry is the contribution.

\subsection{Why H3 Fails: Capacity Is a Proxy, Not the Cause}

The pre-registration assumed firing fraction would scale with
capacity. The actual driver is the magnitude of the per-window
calibration-target shift, which depends on $K$ \textit{and} the
specific window's compositional dynamics. At $K = 100$, the
window-to-window shifts happen to be modest at W3, W4, W6;
the calibration-target trajectory is locally smooth there. At
$K = 200$, every window after W2 has a meaningful shift because
remediation depletes more of the high-priority tail per window.
Capacity is a reasonable proxy for shift magnitude but not a clean
determinant.

\subsection{H2 as the Principled Cost of H1}

The H2 rejection should not be read as a failure of the procedure.
A non-trivial adaptive procedure must, by definition, sometimes
revert to the safe baseline when the recalibration would have
helped. The cost of H1 (no hazard at high $K$) and the cost of H2
(some lost gain at moderate $K$) are not separable; they are
parameterised by the threshold $\tau$. A smaller $\tau$ would fire
less often, reduce the H2 penalty at $K \in \{50, 100\}$, and
re-introduce part of the H1 hazard at $K = 200$. A larger $\tau$
would fire more often and amplify the H2 penalty without further
helping H1.

The pre-registered $\tau = 0.05$ was chosen as the smallest round
threshold below which calibration-target shifts can be argued to be
noise. The H1 success and H2 cost at this threshold are the result
the data deliver.

\subsection{Adaptive vs Static Gate}

H4 is the most operationally interesting result. The static
capacity-conditional gate ($K \leq 100 \to \textsf{lag1}$, $K \geq
200 \to \textsf{fixed}$) requires no per-window decision and no
detector. The adaptive procedure beats it at $K = 200$ by only
$+0.7$~pp, which is within the bootstrap confidence band at most
windows.

A practitioner who knows their fleet's per-window capacity is stable
and falls below $K = 100$ can deploy lag-1 directly without any
detector; if above $K \approx 100$ they can deploy fixed weights
without any detector. The adaptive procedure adds value primarily
when capacity is variable across deployment cycles --- for example,
a team whose patching capacity drops during change-window-restricted
periods and rises during regular operations.

\subsection{What This Result Closes and Opens}

\textit{Closes:} the question of whether \textit{simple} adaptive
switching can solve the Paper~7 + Paper~8 hazard. The answer is
qualified-yes: H1 is supported, but at the cost of H2 and with only
marginal improvement over a static gate.

\textit{Opens:}
\begin{itemize}
\item \textbf{Threshold sensitivity.} A sweep over $\tau$ would map
the H1/H2 trade-off curve; the $\tau = 0.05$ point is one operating
point on a Pareto frontier.
\item \textbf{CUSUM and Bayesian detectors.} A CUSUM
detector~\cite{page1954cusum} that accumulates evidence across
windows might fire more selectively than the one-threshold magnitude
test, reducing the H2 cost.
\item \textbf{Capacity-aware detector.} A detector that conditions on
$K$ explicitly --- e.g., separate $\tau$ per capacity --- could
recover H3 monotonicity by construction at the cost of additional
pre-registration discipline.
\end{itemize}

These are deferred. The contribution of this paper is the negative
boundary: the \textit{simplest} adaptive detector solves H1 but
incurs H2 and produces a non-monotone firing pattern.
